\documentclass[11pt]{article}
\usepackage[utf8]{inputenc}
\usepackage[T1]{fontenc}
\usepackage[spanish]{babel}
\usepackage{amsmath, amssymb, amsthm, mathtools}
\usepackage{braket}
\usepackage{dsfont}
\usepackage[margin=2.5cm]{geometry}
\usepackage{hyperref}
\hypersetup{hidelinks}
\emergencystretch=2em
\allowdisplaybreaks

\newcommand{\ii}{\mathrm{i}}
\newcommand{\ee}{\mathrm{e}}
\newcommand{\adag}{\hat a^\dagger}
\newcommand{\bdag}{\hat b^\dagger}
\newcommand{\Id}{\mathds{1}}
\newcommand{\Omeff}{\Omega_{nm}^{\mathrm{ef}}}

\title{Din\'amica del modelo HC}
\author{Ricardo Del Rosal Gardu\~no \\ \small Doctorado en F\'isica, BUAP}
\date{Notas de doctorado}

\begin{document}
\maketitle

\tableofcontents
\bigskip

\section{Objetivo y punto de partida}

La nota sobre el modelo de Holstein-Cummings construy\'o el hamiltoniano de una
mol\'ecula con una transici\'on electr\'onica, un modo cuantizado de cavidad y
un modo vibracional intramolecular. Despu\'es de trabajar en resonancia
electr\'on-fot\'on, introducir los polaritones y aplicar una aproximaci\'on de
onda rotante al intercambio polarit\'on-fon\'on, el hamiltoniano efectivo del
bloque con $n>0$ excitaciones \'opticas qued\'o escrito como
\begin{equation}
  \boxed{
  \tilde H_{\mathrm{ef}}^{(n)}
  =
  \hbar
  \left[
  \Omega n+\frac{\nu\lambda^2}{4}
  +
  \nu\bdag\hat b
  \right]\Id
  +
  \hbar g\sqrt n\,\hat S_z^{(n)}
  -
  \frac{\hbar\nu\lambda}{2}
  \left(
  \hat b\hat S^{(n)\dagger}
  +
  \bdag\hat S^{(n)}
  \right).
  }
\end{equation}
Aqu\'i, $\Omega$ es la frecuencia com\'un del modo de cavidad y de la
transici\'on electr\'onica en resonancia, $\nu$ es la frecuencia vibracional,
$g$ es el acoplamiento electr\'on-fot\'on y $\lambda$ es el par\'ametro de
Huang--Rhys. Los operadores $\hat S^{(n)}$ y $\hat S^{(n)\dagger}$ bajan y
suben, respectivamente, entre los dos polaritones del bloque $n$.

La interacci\'on conserva el n\'umero efectivo polarit\'on-fon\'on
\begin{equation}
  \boxed{
  \hat M^{(n)}
  =
  \bdag\hat b+\hat S^{(n)\dagger}\hat S^{(n)},
  \qquad
  [\tilde H_{\mathrm{ef}}^{(n)},\hat M^{(n)}]=0.
  }
\end{equation}
Por esta raz\'on, para cada $n>0$ y $m>0$ la din\'amica ocurre dentro del
doblete
\begin{equation}
  \boxed{
  \mathcal H_{nm}
  =
  \operatorname{span}
  \left\{
  \ket{+,n,m-1},
  \ket{-,n,m}
  \right\}.
  }
\end{equation}
El estado $\ket{+,n,m-1}$ contiene al polarit\'on superior y $m-1$ fonones,
mientras que $\ket{-,n,m}$ contiene al polarit\'on inferior y $m$ fonones.
Ambos estados tienen el mismo valor $m$ de $\hat M^{(n)}$. La interacci\'on
convierte uno en el otro sin abandonar el doblete.

El objetivo de esta nota es construir la evoluci\'on temporal desde un estado
inicial general, obtener las amplitudes exactas de cada doblete, reunir todos
los sectores en el estado completo y estudiar una aproximaci\'on especialmente
\'util cuando la vibraci\'on comienza en un estado coherente con un n\'umero
medio grande de fonones.

\section{Estado inicial tripartito}

Se considera inicialmente un estado producto entre el campo, la parte
electr\'onica de la mol\'ecula y su vibraci\'on:
\begin{equation}
  \boxed{
  \ket{\Psi(0)}
  =
  \ket{\alpha}\otimes\ket{\phi}\otimes\ket{\beta}.
  }
\end{equation}
Cada factor se escribe en su base natural. Para el campo,
\begin{equation}
  \ket{\alpha}
  =
  \sum_{n=0}^{\infty}c_n\ket n,
  \qquad
  \sum_{n=0}^{\infty}|c_n|^2=1.
\end{equation}
Para la mol\'ecula de dos niveles,
\begin{equation}
  \ket{\phi}
  =
  c_g\ket g+c_e\ket e,
  \qquad
  |c_g|^2+|c_e|^2=1.
\end{equation}
Finalmente, para la vibraci\'on,
\begin{equation}
  \ket{\beta}
  =
  \sum_{m=0}^{\infty}c_m^{(v)}\ket m,
  \qquad
  \sum_{m=0}^{\infty}|c_m^{(v)}|^2=1.
\end{equation}
Se ha escrito $c_m^{(v)}$ para no confundir los coeficientes vibracionales con
los coeficientes fot\'onicos $c_n$.

\subsection{Estado vibracional en el espacio desplazado}

El hamiltoniano efectivo fue obtenido despu\'es de desplazar la coordenada
vibracional. Por ello, para evolucionar con $\tilde H_{\mathrm{ef}}$ tambi\'en
debe expresarse el estado inicial en el espacio desplazado. En la convenci\'on
empleada en el manuscrito de din\'amica se define
\begin{equation}
  \boxed{
  \ket{\tilde\beta}
  =
  \hat D\!\left(\frac{\lambda}{2}\right)\ket{\beta},
  \qquad
  \hat D(\gamma)
  =
  \exp\!\left(\gamma\bdag-\gamma^*\hat b\right).
  }
\end{equation}
Si el estado vibracional inicial es coherente,
$\ket{\beta}=\hat D(\beta)\ket0$, se utiliza la regla de composici\'on
\begin{equation}
  \hat D(\alpha)\hat D(\beta)
  =
  \exp\!\left[
  \frac{\alpha\beta^*-\alpha^*\beta}{2}
  \right]
  \hat D(\alpha+\beta).
\end{equation}
Tomando $\alpha=\lambda/2$,
\begin{align}
  \ket{\tilde\beta}
  &=
  \hat D\!\left(\frac{\lambda}{2}\right)
  \hat D(\beta)\ket0
  \notag\\
  &=
  \exp\!\left[
  \frac{1}{2}
  \left(
  \frac{\lambda}{2}\beta^*
  -
  \frac{\lambda^*}{2}\beta
  \right)
  \right]
  \hat D\!\left(\beta+\frac{\lambda}{2}\right)\ket0.
\end{align}
Por tanto,
\begin{equation}
  \boxed{
  \ket{\tilde\beta}
  =
  \exp\!\left[
  \frac{\lambda\beta^*-\lambda^*\beta}{4}
  \right]
  \ket{\beta+\lambda/2}.
  }
\end{equation}
Para $\lambda\in\mathbb R$, la exponencial anterior es una fase global:
\begin{equation}
  \exp\!\left[
  \frac{\lambda}{4}(\beta^*-\beta)
  \right]
  =
  \exp\!\left[
  -\ii\frac{\lambda}{2}\operatorname{Im}\beta
  \right].
\end{equation}
Si adem\'as $\beta$ es real, esa fase es igual a uno y simplemente
\begin{equation}
  \boxed{
  \ket{\tilde\beta}
  =
  \ket{\beta+\lambda/2}
  \qquad
  (\lambda,\beta\in\mathbb R).
  }
\end{equation}

La nota de construcci\'on del hamiltoniano usa de manera equivalente la
convenci\'on
$\ket{\tilde\Psi}=\hat D^\dagger(\lambda/2)\ket{\Psi}$. En esa convenci\'on,
el estado coherente desplazado tiene amplitud $\beta-\lambda/2$. Ambas
convenciones describen la misma f\'isica si el signo del desplazamiento se
mantiene de forma consistente. En lo que sigue, toda la dependencia del
estado vibracional inicial se encapsula en los coeficientes
\begin{equation}
  \boxed{
  \ket{\tilde\beta}
  =
  \sum_{m=0}^{\infty}\tilde c_m\ket m,
  \qquad
  \sum_{m=0}^{\infty}|\tilde c_m|^2=1.
  }
\end{equation}
De este modo, las expresiones de la din\'amica no dependen de ocultar una
elecci\'on de signo.

El estado inicial que debe evolucionarse con el hamiltoniano efectivo es
\begin{equation}
  \boxed{
  \ket{\tilde\Psi(0)}
  =
  \left(\sum_{n=0}^{\infty}c_n\ket n\right)
  \left(c_g\ket g+c_e\ket e\right)
  \left(\sum_{m=0}^{\infty}\tilde c_m\ket m\right).
  }
\end{equation}

\section{Parte \'optico-electr\'onica en la base polarit\'onica}

Antes de incorporar la vibraci\'on, se reorganiza el producto entre el campo y
el estado electr\'onico:
\begin{align}
  \ket{\Psi_{\mathrm{op-el}}(0)}
  &=
  \left(\sum_{n=0}^{\infty}c_n\ket n\right)
  \left(c_g\ket g+c_e\ket e\right)
  \notag\\
  &=
  \sum_{n=0}^{\infty}c_nc_g\ket{g,n}
  +
  \sum_{n=0}^{\infty}c_nc_e\ket{e,n}.
\end{align}
Los polaritones del bloque con $n$ excitaciones \'opticas se construyen a
partir de $\ket{e,n-1}$ y $\ket{g,n}$. Por ello, la segunda suma se reindexa
para que ambos estados lleven la misma etiqueta $n$. Sustituyendo
$n\mapsto n-1$,
\begin{equation}
  \sum_{n=0}^{\infty}c_nc_e\ket{e,n}
  =
  \sum_{n=1}^{\infty}c_{n-1}c_e\ket{e,n-1}.
\end{equation}
La suma asociada al estado base contiene adem\'as el estado especial
$\ket{g,0}$, que no tiene pareja electr\'onica dentro del sector de cero
excitaciones. Separ\'andolo,
\begin{equation}
  \sum_{n=0}^{\infty}c_nc_g\ket{g,n}
  =
  c_0c_g\ket{g,0}
  +
  \sum_{n=1}^{\infty}c_nc_g\ket{g,n}.
\end{equation}
Por tanto,
\begin{equation}
  \boxed{
  \ket{\Psi_{\mathrm{op-el}}(0)}
  =
  c_0c_g\ket{g,0}
  +
  \sum_{n=1}^{\infty}
  \left(
  c_{n-1}c_e\ket{e,n-1}
  +
  c_nc_g\ket{g,n}
  \right).
  }
\end{equation}

Para $n>0$, los polaritones resonantes son
\begin{equation}
  \boxed{
  \ket{\pm,n}
  =
  \frac{1}{\sqrt2}
  \left(
  \ket{e,n-1}\pm\ket{g,n}
  \right).
  }
\end{equation}
Sumando y restando estas definiciones se obtienen las relaciones inversas:
\begin{equation}
  \ket{e,n-1}
  =
  \frac{1}{\sqrt2}
  \left(
  \ket{+,n}+\ket{-,n}
  \right),
\end{equation}
\begin{equation}
  \ket{g,n}
  =
  \frac{1}{\sqrt2}
  \left(
  \ket{+,n}-\ket{-,n}
  \right).
\end{equation}
Sustituyendo dentro de cada bloque,
\begin{align}
  &c_{n-1}c_e\ket{e,n-1}
  +
  c_nc_g\ket{g,n}
  \notag\\
  &\quad=
  \frac{c_{n-1}c_e}{\sqrt2}
  \left(\ket{+,n}+\ket{-,n}\right)
  +
  \frac{c_nc_g}{\sqrt2}
  \left(\ket{+,n}-\ket{-,n}\right)
  \notag\\
  &\quad=
  \frac{c_{n-1}c_e+c_nc_g}{\sqrt2}\ket{+,n}
  +
  \frac{c_{n-1}c_e-c_nc_g}{\sqrt2}\ket{-,n}.
\end{align}
Se definen entonces las amplitudes polarit\'onicas iniciales
\begin{equation}
  \boxed{
  d_n^\pm
  =
  \frac{1}{\sqrt2}
  \left(
  c_{n-1}c_e\pm c_nc_g
  \right),
  \qquad n\geq1.
  }
\end{equation}
Con esta notaci\'on, la parte \'optico-electr\'onica queda
\begin{equation}
  \boxed{
  \ket{\Psi_{\mathrm{op-el}}(0)}
  =
  c_0c_g\ket{g,0}
  +
  \sum_{n=1}^{\infty}
  \left(
  d_n^+\ket{+,n}
  +
  d_n^-\ket{-,n}
  \right).
  }
\end{equation}
La amplitud $d_n^+$ mide cu\'anto del estado inicial ocupa el polarit\'on
sim\'etrico, mientras que $d_n^-$ mide la ocupaci\'on del polarit\'on
antisim\'etrico. Estas amplitudes ya contienen tanto la distribuci\'on
fot\'onica como la superposici\'on electr\'onica inicial.

\section{Organizaci\'on del estado inicial en dobletes din\'amicos}

Al recuperar la vibraci\'on desplazada, el estado inicial completo toma la
forma
\begin{align}
  \ket{\tilde\Psi(0)}
  &=
  c_0c_g
  \sum_{m=0}^{\infty}
  \tilde c_m\ket{g,0,m}
  \notag\\
  &\quad+
  \sum_{n=1}^{\infty}\sum_{m=0}^{\infty}
  \tilde c_m
  \left(
  d_n^+\ket{+,n,m}
  +
  d_n^-\ket{-,n,m}
  \right).
\end{align}
Sin embargo, el hamiltoniano efectivo no acopla
$\ket{+,n,m}$ con $\ket{-,n,m}$. Acopla
$\ket{+,n,m-1}$ con $\ket{-,n,m}$. Para exhibir esta estructura se reindexa la
contribuci\'on del polarit\'on superior:
\begin{equation}
  \sum_{m=0}^{\infty}
  \tilde c_md_n^+\ket{+,n,m}
  =
  \sum_{m=1}^{\infty}
  \tilde c_{m-1}d_n^+\ket{+,n,m-1}.
\end{equation}
En el lado derecho, el nuevo \'indice comienza en $m=1$ porque el estado
original con cero fonones corresponde ahora a $m-1=0$.

Para el polarit\'on inferior se separa el t\'ermino $m=0$:
\begin{equation}
  \sum_{m=0}^{\infty}
  \tilde c_md_n^-\ket{-,n,m}
  =
  \tilde c_0d_n^-\ket{-,n,0}
  +
  \sum_{m=1}^{\infty}
  \tilde c_md_n^-\ket{-,n,m}.
\end{equation}
El estado $\ket{-,n,0}$ es aislado: no puede transformarse en
$\ket{+,n,-1}$ porque ese estado no existe. Reuniendo los t\'erminos,
\begin{align}
  \ket{\tilde\Psi(0)}
  &=
  c_0c_g
  \sum_{m=0}^{\infty}
  \tilde c_m\ket{g,0,m}
  +
  \sum_{n=1}^{\infty}
  \tilde c_0d_n^-\ket{-,n,0}
  \notag\\
  &\quad+
  \sum_{n=1}^{\infty}\sum_{m=1}^{\infty}
  \left(
  \tilde c_{m-1}d_n^+\ket{+,n,m-1}
  +
  \tilde c_md_n^-\ket{-,n,m}
  \right).
\end{align}
Los dos primeros t\'erminos son sectores desacoplados. Toda la transferencia
coherente polarit\'on-fon\'on ocurre en la tercera contribuci\'on:
\begin{equation}
  \boxed{
  \ket{\tilde\Psi_{\mathrm{din}}(0)}
  =
  \sum_{n=1}^{\infty}\sum_{m=1}^{\infty}
  \left(
  a_{nm}\ket{+,n,m-1}
  +
  b_{nm}\ket{-,n,m}
  \right),
  }
\end{equation}
donde se han definido, de manera expl\'icita,
\begin{equation}
  \boxed{
  a_{nm}
  =
  \tilde c_{m-1}d_n^+,
  \qquad
  b_{nm}
  =
  \tilde c_md_n^-.
  }
\end{equation}
Para un doblete fijo $(n,m)$,
\begin{equation}
  \boxed{
  \ket{\psi_{nm}(0)}
  =
  a_{nm}\ket{+,n,m-1}
  +
  b_{nm}\ket{-,n,m}.
  }
\end{equation}
La amplitud $a_{nm}$ combina la probabilidad inicial de encontrar el
polarit\'on superior en el bloque $n$ con la amplitud vibracional de tener
$m-1$ fonones. De manera an\'aloga, $b_{nm}$ combina el polarit\'on inferior
con $m$ fonones.

\section{Eigenestados de cada doblete}

La diagonalizaci\'on del hamiltoniano efectivo produce, para cada $n>0$ y
$m>0$, las cantidades
\begin{equation}
  \boxed{
  \begin{aligned}
  A_{nm}
  &=
  \Omega n+\frac{\nu\lambda^2}{4}
  +
  \nu\left(m-\frac12\right),
  \\
  \delta_n
  &=
  2g\sqrt n-\nu,
  \\
  k_m
  &=
  \nu\lambda\sqrt m,
  \\
  \Omega_{nm}^{\mathrm{ef}}
  &=
  \sqrt{\delta_n^2+k_m^2}
  =
  \sqrt{(2g\sqrt n-\nu)^2+\nu^2\lambda^2m}.
  \end{aligned}
  }
\end{equation}
La frecuencia $A_{nm}$ es el centro energ\'etico com\'un del doblete.
La desintonizaci\'on $\delta_n$ compara la separaci\'on polarit\'onica
$2g\sqrt n$ con la frecuencia vibracional $\nu$. La cantidad $k_m$ es la
intensidad del acoplamiento dentro del doblete y crece como $\sqrt m$, como es
caracter\'istico de un operador bos\'onico. Finalmente,
$\Omega_{nm}^{\mathrm{ef}}$ es la frecuencia de Rabi efectiva que combina
desintonizaci\'on y acoplamiento.

Los coeficientes del \'angulo de mezcla son
\begin{equation}
  \boxed{
  C_{nm}
  =
  \sqrt{
  \frac{\Omega_{nm}^{\mathrm{ef}}+\delta_n}
  {2\Omega_{nm}^{\mathrm{ef}}}
  }
  =
  \cos\frac{\theta_{nm}}{2},
  \qquad
  S_{nm}
  =
  \sqrt{
  \frac{\Omega_{nm}^{\mathrm{ef}}-\delta_n}
  {2\Omega_{nm}^{\mathrm{ef}}}
  }
  =
  \sin\frac{\theta_{nm}}{2}.
  }
\end{equation}
De estas definiciones se siguen inmediatamente las identidades
\begin{equation}
  C_{nm}^2+S_{nm}^2=1,
\end{equation}
\begin{equation}
  \boxed{
  C_{nm}^2-S_{nm}^2
  =
  \frac{\delta_n}{\Omega_{nm}^{\mathrm{ef}}},
  \qquad
  2C_{nm}S_{nm}
  =
  \frac{k_m}{\Omega_{nm}^{\mathrm{ef}}}.
  }
\end{equation}
Los eigenestados normalizados son
\begin{equation}
  \boxed{
  \ket{E_{nm}^{+}}
  =
  C_{nm}\ket{+,n,m-1}
  -
  S_{nm}\ket{-,n,m},
  }
\end{equation}
\begin{equation}
  \boxed{
  \ket{E_{nm}^{-}}
  =
  S_{nm}\ket{+,n,m-1}
  +
  C_{nm}\ket{-,n,m}.
  }
\end{equation}
Sus energ\'ias son
\begin{equation}
  \boxed{
  E_{nm}^{\pm}
  =
  \hbar
  \left(
  A_{nm}
  \pm
  \frac{\Omega_{nm}^{\mathrm{ef}}}{2}
  \right).
  }
\end{equation}

Para obtener las relaciones inversas se escriben las dos ecuaciones anteriores
en forma matricial:
\begin{equation}
  \begin{pmatrix}
  \ket{E_{nm}^{+}}\\
  \ket{E_{nm}^{-}}
  \end{pmatrix}
  =
  \begin{pmatrix}
  C_{nm} & -S_{nm}\\
  S_{nm} & C_{nm}
  \end{pmatrix}
  \begin{pmatrix}
  \ket{+,n,m-1}\\
  \ket{-,n,m}
  \end{pmatrix}.
\end{equation}
La matriz de mezcla es ortogonal, por lo que su inversa es su transpuesta.
Entonces,
\begin{equation}
  \boxed{
  \ket{+,n,m-1}
  =
  C_{nm}\ket{E_{nm}^{+}}
  +
  S_{nm}\ket{E_{nm}^{-}},
  }
\end{equation}
\begin{equation}
  \boxed{
  \ket{-,n,m}
  =
  -S_{nm}\ket{E_{nm}^{+}}
  +
  C_{nm}\ket{E_{nm}^{-}}.
  }
\end{equation}

\section{Evoluci\'on exacta de un doblete}

Sustituyendo las relaciones inversas en el estado inicial del doblete,
\begin{align}
  \ket{\psi_{nm}(0)}
  &=
  a_{nm}
  \left(
  C_{nm}\ket{E_{nm}^{+}}
  +
  S_{nm}\ket{E_{nm}^{-}}
  \right)
  \notag\\
  &\quad+
  b_{nm}
  \left(
  -S_{nm}\ket{E_{nm}^{+}}
  +
  C_{nm}\ket{E_{nm}^{-}}
  \right).
\end{align}
Agrupando los coeficientes de cada eigenestado,
\begin{equation}
  \boxed{
  \ket{\psi_{nm}(0)}
  =
  \left(a_{nm}C_{nm}-b_{nm}S_{nm}\right)\ket{E_{nm}^{+}}
  +
  \left(a_{nm}S_{nm}+b_{nm}C_{nm}\right)\ket{E_{nm}^{-}}.
  }
\end{equation}

El operador de evoluci\'on del doblete es
\begin{equation}
  \hat U_{nm}(t)
  =
  \exp\!\left[
  -\frac{\ii}{\hbar}
  \tilde H_{\mathrm{ef}}^{(n,m)}t
  \right].
\end{equation}
Como $\ket{E_{nm}^{\pm}}$ son eigenestados,
\begin{equation}
  \hat U_{nm}(t)\ket{E_{nm}^{\pm}}
  =
  \exp\!\left(
  -\frac{\ii}{\hbar}E_{nm}^{\pm}t
  \right)
  \ket{E_{nm}^{\pm}}.
\end{equation}
Por ello,
\begin{align}
  \ket{\psi_{nm}(t)}
  &=
  \ee^{-\ii(A_{nm}+\Omega_{nm}^{\mathrm{ef}}/2)t}
  \left(a_{nm}C_{nm}-b_{nm}S_{nm}\right)
  \ket{E_{nm}^{+}}
  \notag\\
  &\quad+
  \ee^{-\ii(A_{nm}-\Omega_{nm}^{\mathrm{ef}}/2)t}
  \left(a_{nm}S_{nm}+b_{nm}C_{nm}\right)
  \ket{E_{nm}^{-}}.
\end{align}
El factor $\ee^{-\ii A_{nm}t}$ es com\'un a las dos ramas y representa la fase
debida al centro energ\'etico del doblete:
\begin{align}
  \ket{\psi_{nm}(t)}
  &=
  \ee^{-\ii A_{nm}t}
  \Big[
  \ee^{-\ii\Omega_{nm}^{\mathrm{ef}}t/2}
  \left(a_{nm}C_{nm}-b_{nm}S_{nm}\right)
  \ket{E_{nm}^{+}}
  \notag\\
  &\hspace{3.3cm}+
  \ee^{+\ii\Omega_{nm}^{\mathrm{ef}}t/2}
  \left(a_{nm}S_{nm}+b_{nm}C_{nm}\right)
  \ket{E_{nm}^{-}}
  \Big].
\end{align}

\subsection{Regreso a la base polarit\'on-fon\'on}

Se define
\begin{equation}
  \boxed{
  \ket{\psi_{nm}(t)}
  =
  p_{nm}^{+}(t)\ket{+,n,m-1}
  +
  p_{nm}^{-}(t)\ket{-,n,m}.
  }
\end{equation}
Para calcular $p_{nm}^{+}(t)$ y $p_{nm}^{-}(t)$ se sustituyen nuevamente los
eigenestados en t\'erminos de la base polarit\'on-fon\'on:
\begin{align}
  p_{nm}^{+}(t)
  &=
  \ee^{-\ii A_{nm}t}
  \Big[
  C_{nm}\ee^{-\ii\Omega_{nm}^{\mathrm{ef}}t/2}
  \left(a_{nm}C_{nm}-b_{nm}S_{nm}\right)
  \notag\\
  &\hspace{3.2cm}+
  S_{nm}\ee^{+\ii\Omega_{nm}^{\mathrm{ef}}t/2}
  \left(a_{nm}S_{nm}+b_{nm}C_{nm}\right)
  \Big],
\end{align}
\begin{align}
  p_{nm}^{-}(t)
  &=
  \ee^{-\ii A_{nm}t}
  \Big[
  -S_{nm}\ee^{-\ii\Omega_{nm}^{\mathrm{ef}}t/2}
  \left(a_{nm}C_{nm}-b_{nm}S_{nm}\right)
  \notag\\
  &\hspace{3.2cm}+
  C_{nm}\ee^{+\ii\Omega_{nm}^{\mathrm{ef}}t/2}
  \left(a_{nm}S_{nm}+b_{nm}C_{nm}\right)
  \Big].
\end{align}

Para simplificar la primera amplitud se agrupan por separado los t\'erminos
proporcionales a $a_{nm}$ y a $b_{nm}$:
\begin{align}
  p_{nm}^{+}(t)
  &=
  \ee^{-\ii A_{nm}t}
  \Big[
  a_{nm}
  \left(
  C_{nm}^2\ee^{-\ii x}
  +
  S_{nm}^2\ee^{+\ii x}
  \right)
  \notag\\
  &\hspace{3.2cm}+
  b_{nm}C_{nm}S_{nm}
  \left(
  \ee^{+\ii x}-\ee^{-\ii x}
  \right)
  \Big],
\end{align}
donde se defini\'o temporalmente
\begin{equation}
  x=\frac{\Omega_{nm}^{\mathrm{ef}}t}{2}.
\end{equation}
Usando
\begin{equation}
  \ee^{+\ii x}+\ee^{-\ii x}=2\cos x,
  \qquad
  \ee^{+\ii x}-\ee^{-\ii x}=2\ii\sin x,
\end{equation}
se obtiene
\begin{align}
  C_{nm}^2\ee^{-\ii x}+S_{nm}^2\ee^{+\ii x}
  &=
  (C_{nm}^2+S_{nm}^2)\cos x
  \notag\\
  &\quad
  -\ii(C_{nm}^2-S_{nm}^2)\sin x
  \notag\\
  &=
  \cos x
  -
  \ii\frac{\delta_n}{\Omega_{nm}^{\mathrm{ef}}}\sin x,
\end{align}
y
\begin{equation}
  C_{nm}S_{nm}
  \left(\ee^{+\ii x}-\ee^{-\ii x}\right)
  =
  \ii\frac{k_m}{\Omega_{nm}^{\mathrm{ef}}}\sin x.
\end{equation}
El mismo procedimiento para $p_{nm}^{-}(t)$ produce el resultado exacto
\begin{equation}
  \boxed{
  \begin{aligned}
  p_{nm}^{+}(t)
  &=
  \ee^{-\ii A_{nm}t}
  \Bigg[
  \left(
  \cos\frac{\Omega_{nm}^{\mathrm{ef}}t}{2}
  -
  \ii\frac{\delta_n}{\Omega_{nm}^{\mathrm{ef}}}
  \sin\frac{\Omega_{nm}^{\mathrm{ef}}t}{2}
  \right)a_{nm}
  \\
  &\hspace{2.3cm}+
  \ii\frac{k_m}{\Omega_{nm}^{\mathrm{ef}}}
  \sin\frac{\Omega_{nm}^{\mathrm{ef}}t}{2}
  \,b_{nm}
  \Bigg],
  \\
  p_{nm}^{-}(t)
  &=
  \ee^{-\ii A_{nm}t}
  \Bigg[
  \ii\frac{k_m}{\Omega_{nm}^{\mathrm{ef}}}
  \sin\frac{\Omega_{nm}^{\mathrm{ef}}t}{2}
  \,a_{nm}
  \\
  &\hspace{2.3cm}+
  \left(
  \cos\frac{\Omega_{nm}^{\mathrm{ef}}t}{2}
  +
  \ii\frac{\delta_n}{\Omega_{nm}^{\mathrm{ef}}}
  \sin\frac{\Omega_{nm}^{\mathrm{ef}}t}{2}
  \right)b_{nm}
  \Bigg].
  \end{aligned}
  }
\end{equation}
Al sustituir $a_{nm}$ y $b_{nm}$, las amplitudes quedan directamente en
t\'erminos del estado inicial:
\begin{equation}
  \boxed{
  \begin{aligned}
  p_{nm}^{+}(t)
  &=
  \ee^{-\ii A_{nm}t}
  \Bigg[
  \left(
  \cos\frac{\Omega_{nm}^{\mathrm{ef}}t}{2}
  -
  \ii\frac{\delta_n}{\Omega_{nm}^{\mathrm{ef}}}
  \sin\frac{\Omega_{nm}^{\mathrm{ef}}t}{2}
  \right)
  \tilde c_{m-1}d_n^+
  \\
  &\hspace{2.1cm}+
  \ii\frac{k_m}{\Omega_{nm}^{\mathrm{ef}}}
  \sin\frac{\Omega_{nm}^{\mathrm{ef}}t}{2}
  \tilde c_md_n^-
  \Bigg],
  \\
  p_{nm}^{-}(t)
  &=
  \ee^{-\ii A_{nm}t}
  \Bigg[
  \ii\frac{k_m}{\Omega_{nm}^{\mathrm{ef}}}
  \sin\frac{\Omega_{nm}^{\mathrm{ef}}t}{2}
  \tilde c_{m-1}d_n^+
  \\
  &\hspace{2.1cm}+
  \left(
  \cos\frac{\Omega_{nm}^{\mathrm{ef}}t}{2}
  +
  \ii\frac{\delta_n}{\Omega_{nm}^{\mathrm{ef}}}
  \sin\frac{\Omega_{nm}^{\mathrm{ef}}t}{2}
  \right)
  \tilde c_md_n^-
  \Bigg].
  \end{aligned}
  }
\end{equation}
Estas dos amplitudes muestran las dos partes de la evoluci\'on. Los t\'erminos
con coseno y con $\delta_n/\Omega_{nm}^{\mathrm{ef}}$ describen la amplitud
que permanece en la misma rama polarit\'onica. Los t\'erminos con
$k_m/\Omega_{nm}^{\mathrm{ef}}$ describen la transferencia entre las dos
ramas acompa\~nada por la absorci\'on o emisi\'on de un fon\'on.

\section{Evoluci\'on exacta del estado completo}

La parte din\'amica del estado se obtiene sumando todos los dobletes:
\begin{equation}
  \boxed{
  \ket{\tilde\Psi_{\mathrm{din}}(t)}
  =
  \sum_{n=1}^{\infty}\sum_{m=1}^{\infty}
  \left[
  p_{nm}^{+}(t)\ket{+,n,m-1}
  +
  p_{nm}^{-}(t)\ket{-,n,m}
  \right].
  }
\end{equation}
Sin embargo, esta suma no contiene los sectores desacoplados que estaban
presentes en el estado inicial.

Para $n=0$, los estados $\ket{g,0,m}$ tienen energ\'ia
\begin{equation}
  E_{g,0,m}=\hbar\nu m.
\end{equation}
Por tanto, su evoluci\'on es
\begin{equation}
  c_0c_g
  \sum_{m=0}^{\infty}
  \tilde c_m
  \ee^{-\ii\nu mt}
  \ket{g,0,m}.
\end{equation}
Para $n>0$ y $m=0$, el estado aislado $\ket{-,n,0}$ tiene energ\'ia
\begin{equation}
  E_{n,0}^{-}
  =
  \hbar
  \left(
  \Omega n+\frac{\nu\lambda^2}{4}-g\sqrt n
  \right),
\end{equation}
por lo que solo acumula la fase
$\exp(-\ii E_{n,0}^{-}t/\hbar)$.

Reuniendo los tres tipos de sectores se obtiene el estado completo:
\begin{equation}
  \boxed{
  \begin{aligned}
  \ket{\tilde\Psi(t)}
  &=
  c_0c_g
  \sum_{m=0}^{\infty}
  \tilde c_m
  \ee^{-\ii\nu mt}
  \ket{g,0,m}
  \\
  &\quad+
  \sum_{n=1}^{\infty}
  \tilde c_0d_n^-
  \ee^{-\ii E_{n,0}^{-}t/\hbar}
  \ket{-,n,0}
  \\
  &\quad+
  \sum_{n=1}^{\infty}\sum_{m=1}^{\infty}
  \left[
  p_{nm}^{+}(t)\ket{+,n,m-1}
  +
  p_{nm}^{-}(t)\ket{-,n,m}
  \right].
  \end{aligned}
  }
\end{equation}
Esta es la expresi\'on exacta que conserva toda la informaci\'on del estado
inicial. Los dos primeros renglones evolucionan mediante fases porque no
disponen de un estado compa\~nero permitido. El tercer rengl\'on contiene la
din\'amica no trivial.

\subsection{Diccionario completo de las cantidades del resultado}

Para que la soluci\'on anterior pueda utilizarse sin regresar a secciones
previas, todas sus definiciones se re\'unen a continuaci\'on:
\begin{equation}
  \boxed{
  \begin{aligned}
  d_n^\pm
  &=
  \frac{c_{n-1}c_e\pm c_nc_g}{\sqrt2},
  \\
  a_{nm}
  &=
  \tilde c_{m-1}d_n^+,
  \qquad
  b_{nm}
  =
  \tilde c_md_n^-,
  \\
  A_{nm}
  &=
  \Omega n+\frac{\nu\lambda^2}{4}
  +
  \nu\left(m-\frac12\right),
  \\
  \delta_n
  &=
  2g\sqrt n-\nu,
  \qquad
  k_m
  =
  \nu\lambda\sqrt m,
  \\
  \Omega_{nm}^{\mathrm{ef}}
  &=
  \sqrt{\delta_n^2+k_m^2},
  \\
  C_{nm}
  &=
  \sqrt{
  \frac{\Omega_{nm}^{\mathrm{ef}}+\delta_n}
  {2\Omega_{nm}^{\mathrm{ef}}}
  },
  \qquad
  S_{nm}
  =
  \sqrt{
  \frac{\Omega_{nm}^{\mathrm{ef}}-\delta_n}
  {2\Omega_{nm}^{\mathrm{ef}}}
  },
  \\
  E_{nm}^{\pm}
  &=
  \hbar
  \left(
  A_{nm}\pm\frac{\Omega_{nm}^{\mathrm{ef}}}{2}
  \right).
  \end{aligned}
  }
\end{equation}

\section{Regreso a la base electr\'onica}

Los polaritones mezclan la mol\'ecula y el campo. Para calcular cantidades como
la poblaci\'on electr\'onica es conveniente regresar a la base
$\{\ket g,\ket e\}$. Recordando
\begin{equation}
  \ket{+,n}
  =
  \frac{1}{\sqrt2}
  \left(
  \ket{e,n-1}+\ket{g,n}
  \right),
\end{equation}
\begin{equation}
  \ket{-,n}
  =
  \frac{1}{\sqrt2}
  \left(
  \ket{e,n-1}-\ket{g,n}
  \right),
\end{equation}
al incluir los fonones se tiene
\begin{equation}
  \ket{+,n,m-1}
  =
  \frac{1}{\sqrt2}
  \left(
  \ket{n-1,m-1}\ket e
  +
  \ket{n,m-1}\ket g
  \right),
\end{equation}
\begin{equation}
  \ket{-,n,m}
  =
  \frac{1}{\sqrt2}
  \left(
  \ket{n-1,m}\ket e
  -
  \ket{n,m}\ket g
  \right),
\end{equation}
donde $\ket{n,m}$ representa ahora el estado conjunto de campo y vibraci\'on.

Al sustituir en la parte din\'amica,
\begin{align}
  \ket{\tilde\Psi_{\mathrm{din}}(t)}
  &=
  \frac{1}{\sqrt2}
  \sum_{n=1}^{\infty}\sum_{m=1}^{\infty}
  \Big[
  p_{nm}^{+}(t)
  \left(
  \ket{n-1,m-1}\ket e+\ket{n,m-1}\ket g
  \right)
  \notag\\
  &\hspace{3.4cm}+
  p_{nm}^{-}(t)
  \left(
  \ket{n-1,m}\ket e-\ket{n,m}\ket g
  \right)
  \Big].
\end{align}
Agrupando seg\'un el estado electr\'onico,
\begin{equation}
  \boxed{
  \ket{\tilde\Psi_{\mathrm{din}}(t)}
  =
  \ket{\chi_g(t)}\ket g
  +
  \ket{\chi_e(t)}\ket e,
  }
\end{equation}
con
\begin{equation}
  \boxed{
  \ket{\chi_g(t)}
  =
  \frac{1}{\sqrt2}
  \sum_{n=1}^{\infty}\sum_{m=1}^{\infty}
  \left[
  p_{nm}^{+}(t)\ket{n,m-1}
  -
  p_{nm}^{-}(t)\ket{n,m}
  \right],
  }
\end{equation}
\begin{equation}
  \boxed{
  \ket{\chi_e(t)}
  =
  \frac{1}{\sqrt2}
  \sum_{n=1}^{\infty}\sum_{m=1}^{\infty}
  \left[
  p_{nm}^{+}(t)\ket{n-1,m-1}
  +
  p_{nm}^{-}(t)\ket{n-1,m}
  \right].
  }
\end{equation}
Los estados $\ket{\chi_g(t)}$ y $\ket{\chi_e(t)}$ no son, en general,
estados normalizados. Sus normas son las probabilidades de encontrar a la
mol\'ecula en $\ket g$ o $\ket e$ dentro del sector din\'amico. Adem\'as,
cada uno conserva toda la correlaci\'on residual entre el campo y la
vibraci\'on condicionada al resultado electr\'onico.

\section{Vibraci\'on inicial coherente}

Se considera ahora el caso en que el estado vibracional desplazado es
coherente:
\begin{equation}
  \boxed{
  \ket{\tilde\beta}
  =
  \sum_{m=0}^{\infty}
  \tilde c_m\ket m,
  \qquad
  \tilde c_m
  =
  \ee^{-|\tilde\beta|^2/2}
  \frac{\tilde\beta^m}{\sqrt{m!}}.
  }
\end{equation}
La distribuci\'on de fonones es de Poisson:
\begin{equation}
  P_m
  =
  |\tilde c_m|^2
  =
  \ee^{-M}\frac{M^m}{m!},
  \qquad
  \boxed{M=|\tilde\beta|^2}.
\end{equation}
Su media y su varianza son
\begin{equation}
  \langle m\rangle=M,
  \qquad
  (\Delta m)^2=M,
  \qquad
  \Delta m=\sqrt M.
\end{equation}
Cuando $M\gg1$, la distribuci\'on est\'a concentrada alrededor de $m=M$ y su
anchura relativa es peque\~na:
\begin{equation}
  \frac{\Delta m}{M}
  =
  \frac{1}{\sqrt M}
  \ll1.
\end{equation}
Esto permite aproximar las cantidades que cambian lentamente con $m$ por su
valor cerca del centro de la distribuci\'on.

\subsection{Relaci\'on entre coeficientes vecinos}

Para un estado coherente,
\begin{align}
  \frac{\tilde c_m}{\tilde c_{m-1}}
  &=
  \frac{
  \ee^{-|\tilde\beta|^2/2}
  \tilde\beta^m/\sqrt{m!}
  }{
  \ee^{-|\tilde\beta|^2/2}
  \tilde\beta^{m-1}/\sqrt{(m-1)!}
  }
  \notag\\
  &=
  \tilde\beta
  \sqrt{\frac{(m-1)!}{m!}}
  =
  \frac{\tilde\beta}{\sqrt m}.
\end{align}
Para los valores dominantes $m\simeq M=|\tilde\beta|^2$,
\begin{equation}
  \frac{\tilde c_m}{\tilde c_{m-1}}
  \simeq
  \frac{\tilde\beta}{|\tilde\beta|}.
\end{equation}
Si $\tilde\beta$ es real y positivo, la fase anterior es uno y
\begin{equation}
  \boxed{
  \tilde c_m\simeq\tilde c_{m-1}
  \qquad
  (m\simeq M,\ \tilde\beta>0).
  }
\end{equation}
Esta aproximaci\'on permite reemplazar las amplitudes vibracionales de los dos
estados de cada doblete por un mismo coeficiente local. Si $\tilde\beta$ es
complejo, debe conservarse la fase
$\tilde\beta/|\tilde\beta|$.

\section{Linealizaci\'on de la frecuencia efectiva}

La frecuencia de Rabi efectiva depende de $m$ como
\begin{equation}
  \Omega_{nm}^{\mathrm{ef}}
  =
  \sqrt{\delta_n^2+\nu^2\lambda^2m}.
\end{equation}
Para linealizarla alrededor de $m=M$ se escribe su expansi\'on de Taylor de
primer orden:
\begin{equation}
  \Omega_{nm}^{\mathrm{ef}}
  \simeq
  \Omega_{nM}^{\mathrm{ef}}
  +
  (m-M)
  \left.
  \frac{\partial\Omega_{nm}^{\mathrm{ef}}}{\partial m}
  \right|_{m=M}.
\end{equation}
La derivada se calcula paso a paso:
\begin{align}
  \frac{\partial\Omega_{nm}^{\mathrm{ef}}}{\partial m}
  &=
  \frac{\partial}{\partial m}
  \left(\delta_n^2+\nu^2\lambda^2m\right)^{1/2}
  \notag\\
  &=
  \frac12
  \left(\delta_n^2+\nu^2\lambda^2m\right)^{-1/2}
  \nu^2\lambda^2
  \notag\\
  &=
  \frac{\nu^2\lambda^2}
  {2\Omega_{nm}^{\mathrm{ef}}}.
\end{align}
Se define
\begin{equation}
  \boxed{
  \Omega_{nM}^{\mathrm{ef}\prime}
  =
  \left.
  \frac{\partial\Omega_{nm}^{\mathrm{ef}}}{\partial m}
  \right|_{m=M}
  =
  \frac{\nu^2\lambda^2}
  {2\Omega_{nM}^{\mathrm{ef}}}.
  }
\end{equation}
Entonces,
\begin{align}
  \Omega_{nm}^{\mathrm{ef}}
  &\simeq
  \Omega_{nM}^{\mathrm{ef}}
  +(m-M)\Omega_{nM}^{\mathrm{ef}\prime}
  \notag\\
  &=
  \left(
  \Omega_{nM}^{\mathrm{ef}}
  -
  M\Omega_{nM}^{\mathrm{ef}\prime}
  \right)
  +
  m\Omega_{nM}^{\mathrm{ef}\prime}.
\end{align}
Definiendo
\begin{equation}
  \boxed{
  \varphi_{nM}
  =
  \Omega_{nM}^{\mathrm{ef}}
  -
  M\Omega_{nM}^{\mathrm{ef}\prime},
  }
\end{equation}
la linealizaci\'on adopta la forma
\begin{equation}
  \boxed{
  \Omega_{nm}^{\mathrm{ef}}
  \simeq
  \varphi_{nM}
  +
  m\Omega_{nM}^{\mathrm{ef}\prime}.
  }
\end{equation}
La primera cantidad produce una fase global para cada rama del doblete,
mientras que el t\'ermino lineal en $m$ rota la amplitud del estado coherente
en el espacio de fase.

En la misma regi\'on estrecha de valores de $m$, los eigenvectores var\'ian
lentamente y se aproximan por
\begin{equation}
  \boxed{
  C_{nm}\simeq C_{nM}\equiv C,
  \qquad
  S_{nm}\simeq S_{nM}\equiv S.
  }
\end{equation}
La notaci\'on abreviada $C$ y $S$ debe entenderse dentro de un bloque $n$
fijo; sus valores todav\'ia pueden depender de $n$ y de la media $M$.

\section{Evoluci\'on aproximada para un paquete coherente}

Usando $\tilde c_{m-1}\simeq\tilde c_m$, el estado inicial de un doblete se
aproxima por
\begin{equation}
  \ket{\psi_{nm}(0)}
  \simeq
  \tilde c_m
  \left(
  d_n^+\ket{+,n,m-1}
  +
  d_n^-\ket{-,n,m}
  \right).
\end{equation}
En la base de energ\'ia,
\begin{equation}
  \boxed{
  \ket{\psi_{nm}(0)}
  \simeq
  \tilde c_m
  \left(d_n^+C-d_n^-S\right)\ket{E_{nm}^{+}}
  +
  \tilde c_m
  \left(d_n^+S+d_n^-C\right)\ket{E_{nm}^{-}}.
  }
\end{equation}
Despu\'es de evolucionar,
\begin{align}
  \ket{\psi_{nm}(t)}
  &\simeq
  \tilde c_m\ee^{-\ii A_{nM}t}
  \Big[
  \left(d_n^+C-d_n^-S\right)
  \ee^{-\ii\Omega_{nm}^{\mathrm{ef}}t/2}
  \ket{E_{nm}^{+}}
  \notag\\
  &\hspace{3.1cm}+
  \left(d_n^+S+d_n^-C\right)
  \ee^{+\ii\Omega_{nm}^{\mathrm{ef}}t/2}
  \ket{E_{nm}^{-}}
  \Big],
\end{align}
donde
\begin{equation}
  \boxed{
  A_{nM}
  =
  \Omega n+\frac{\nu\lambda^2}{4}
  +
  \nu\left(M-\frac12\right).
  }
\end{equation}
Al usar la frecuencia linealizada,
\begin{equation}
  \ee^{-\ii\Omega_{nm}^{\mathrm{ef}}t/2}
  \simeq
  \ee^{-\ii\varphi_{nM}t/2}
  \ee^{-\ii m\Omega_{nM}^{\mathrm{ef}\prime}t/2},
\end{equation}
\begin{equation}
  \ee^{+\ii\Omega_{nm}^{\mathrm{ef}}t/2}
  \simeq
  \ee^{+\ii\varphi_{nM}t/2}
  \ee^{+\ii m\Omega_{nM}^{\mathrm{ef}\prime}t/2}.
\end{equation}

\subsection{Conversi\'on de las sumas en estados coherentes rotados}

Para cualquier \'angulo $\gamma$,
\begin{align}
  \sum_{m=0}^{\infty}
  \tilde c_m\ee^{-\ii m\gamma}\ket m
  &=
  \ee^{-|\tilde\beta|^2/2}
  \sum_{m=0}^{\infty}
  \frac{
  \left(\tilde\beta\ee^{-\ii\gamma}\right)^m
  }{\sqrt{m!}}
  \ket m
  \notag\\
  &=
  \ket{\tilde\beta\ee^{-\ii\gamma}}.
\end{align}
Por tanto,
\begin{equation}
  \boxed{
  \sum_{m=0}^{\infty}
  \tilde c_m\ee^{\mp\ii m\gamma}\ket m
  =
  \ket{\tilde\beta\ee^{\mp\ii\gamma}}.
  }
\end{equation}
Para una suma que contiene $\ket{m-1}$, se cambia la variable
$r=m-1$:
\begin{align}
  \sum_{m=1}^{\infty}
  \tilde c_m\ee^{-\ii m\gamma}\ket{m-1}
  &=
  \sum_{r=0}^{\infty}
  \tilde c_{r+1}\ee^{-\ii(r+1)\gamma}\ket r
  \notag\\
  &\simeq
  \ee^{-\ii\gamma}
  \sum_{r=0}^{\infty}
  \tilde c_r\ee^{-\ii r\gamma}\ket r
  \notag\\
  &=
  \ee^{-\ii\gamma}
  \ket{\tilde\beta\ee^{-\ii\gamma}}.
\end{align}
An\'alogamente,
\begin{equation}
  \boxed{
  \sum_{m=1}^{\infty}
  \tilde c_m\ee^{\mp\ii m\gamma}\ket{m-1}
  \simeq
  \ee^{\mp\ii\gamma}
  \ket{\tilde\beta\ee^{\mp\ii\gamma}}.
  }
\end{equation}
La fase adicional aparece porque el estado vibracional del polarit\'on
superior contiene un fon\'on menos.

\subsection{Estado aproximado completo del sector din\'amico}

Tomando
\begin{equation}
  \gamma_{nM}(t)
  =
  \frac{\Omega_{nM}^{\mathrm{ef}\prime}t}{2},
\end{equation}
y sumando sobre $m$, se obtiene
\begin{equation}
  \boxed{
  \begin{aligned}
  \ket{\tilde\Psi_{\mathrm{din}}(t)}
  &\simeq
  \sum_{n=1}^{\infty}
  \ee^{-\ii A_{nM}t}
  \Bigg\{
  \left(d_n^+C-d_n^-S\right)
  \ee^{-\ii\varphi_{nM}t/2}
  \ket{\tilde\beta\ee^{-\ii\gamma_{nM}(t)}}
  \\
  &\qquad\qquad\otimes
  \left[
  C\ee^{-\ii\gamma_{nM}(t)}\ket{+,n}
  -
  S\ket{-,n}
  \right]
  \\
  &\quad+
  \left(d_n^+S+d_n^-C\right)
  \ee^{+\ii\varphi_{nM}t/2}
  \ket{\tilde\beta\ee^{+\ii\gamma_{nM}(t)}}
  \\
  &\qquad\qquad\otimes
  \left[
  S\ee^{+\ii\gamma_{nM}(t)}\ket{+,n}
  +
  C\ket{-,n}
  \right]
  \Bigg\}.
  \end{aligned}
  }
\end{equation}
El resultado muestra una separaci\'on f\'isica muy clara. El paquete
vibracional coherente inicial se divide en dos componentes que giran en
sentidos opuestos en el espacio de fase, con amplitudes
$\tilde\beta\exp[\mp\ii\gamma_{nM}(t)]$. Cada componente queda correlacionada
con una superposici\'on polarit\'onica distinta. La din\'amica efectiva crea,
por tanto, entrelazamiento entre polaritones y vibraci\'on.

La separaci\'on angular entre los dos paquetes vibracionales es
\begin{equation}
  \boxed{
  \Delta\vartheta_{nM}(t)
  =
  2\gamma_{nM}(t)
  =
  \Omega_{nM}^{\mathrm{ef}\prime}t.
  }
\end{equation}
Mientras los dos paquetes permanecen superpuestos, las ramas polarit\'onicas
pueden interferir. Cuando se separan en el espacio de fase, la vibraci\'on
almacena informaci\'on sobre la rama polarit\'onica y reduce su interferencia.

\section{Variables din\'amicas}

El hamiltoniano efectivo describe un intercambio entre polaritones y fonones,
\begin{equation}
  -\frac{\hbar\nu\lambda}{2}
  \left(
  \hat b\hat S^{(n)\dagger}
  +
  \bdag\hat S^{(n)}
  \right).
\end{equation}
Sin embargo, pueden estudiarse dos niveles distintos de din\'amica. En la base
desnuda interesa la transferencia entre $\ket g$ y $\ket e$. En la base
vestida interesa la transferencia entre $\ket{+,n}$ y $\ket{-,n}$.

\subsection{Inversi\'on electr\'onica}

La inversi\'on electr\'onica se define como
\begin{equation}
  \boxed{
  W(t)
  =
  \langle\hat\sigma_z(t)\rangle
  =
  \bra{\tilde\Psi(t)}
  \hat\sigma_z
  \ket{\tilde\Psi(t)}
  =
  P_e(t)-P_g(t),
  }
\end{equation}
donde
\begin{equation}
  \hat\sigma_z
  =
  \ket e\bra e-\ket g\bra g.
\end{equation}
Las poblaciones se calculan con los proyectores
\begin{equation}
  \boxed{
  P_e(t)
  =
  \bra{\tilde\Psi(t)}
  \left(
  \ket e\bra e
  \otimes\Id_{\mathrm{cav}}
  \otimes\Id_{\mathrm{vib}}
  \right)
  \ket{\tilde\Psi(t)},
  }
\end{equation}
\begin{equation}
  \boxed{
  P_g(t)
  =
  \bra{\tilde\Psi(t)}
  \left(
  \ket g\bra g
  \otimes\Id_{\mathrm{cav}}
  \otimes\Id_{\mathrm{vib}}
  \right)
  \ket{\tilde\Psi(t)}.
  }
\end{equation}
Si se usa la descomposici\'on
$\ket{\tilde\Psi(t)}=\ket{\chi_g(t)}\ket g+\ket{\chi_e(t)}\ket e$,
incluyendo en $\ket{\chi_g}$ y $\ket{\chi_e}$ los sectores desacoplados
correspondientes, entonces
\begin{equation}
  \boxed{
  W(t)
  =
  \braket{\chi_e(t)|\chi_e(t)}
  -
  \braket{\chi_g(t)|\chi_g(t)}.
  }
\end{equation}
Los valores extremos tienen una interpretaci\'on directa:
\begin{equation}
  W(t)=1
  \quad\Longleftrightarrow\quad
  P_e(t)=1,
\end{equation}
\begin{equation}
  W(t)=-1
  \quad\Longleftrightarrow\quad
  P_g(t)=1,
\end{equation}
\begin{equation}
  W(t)=0
  \quad\Longleftrightarrow\quad
  P_e(t)=P_g(t)=\frac12.
\end{equation}
Aunque $W(t)$ no depende directamente de la coherencia electr\'onica
$\ket e\bra g$, las normas de $\ket{\chi_e(t)}$ y $\ket{\chi_g(t)}$ contienen
la interferencia y las correlaciones creadas con el campo y la vibraci\'on.

\subsection{Inversi\'on polarit\'onica}

Para medir el intercambio que aparece de forma directa en el hamiltoniano
efectivo se define
\begin{equation}
  \hat S_z
  =
  \sum_{n=1}^{\infty}
  \hat S_z^{(n)}
  \otimes\Id_{\mathrm{vib}},
\end{equation}
con
\begin{equation}
  \hat S_z^{(n)}
  =
  \ket{+,n}\bra{+,n}
  -
  \ket{-,n}\bra{-,n}.
\end{equation}
La inversi\'on polarit\'onica es
\begin{equation}
  \boxed{
  \mathcal W(t)
  =
  \langle\hat S_z(t)\rangle
  =
  \bra{\tilde\Psi(t)}
  \hat S_z
  \ket{\tilde\Psi(t)}.
  }
\end{equation}
Definiendo las probabilidades totales de cada polarit\'on, sumadas sobre los
fonones,
\begin{equation}
  P_{+,n}(t)
  =
  \sum_{m=0}^{\infty}
  \left|
  \braket{+,n,m|\tilde\Psi(t)}
  \right|^2,
\end{equation}
\begin{equation}
  P_{-,n}(t)
  =
  \sum_{m=0}^{\infty}
  \left|
  \braket{-,n,m|\tilde\Psi(t)}
  \right|^2,
\end{equation}
se obtiene
\begin{equation}
  \boxed{
  \mathcal W(t)
  =
  \sum_{n=1}^{\infty}
  \left[
  P_{+,n}(t)-P_{-,n}(t)
  \right].
  }
\end{equation}
$\mathcal W(t)$ observa directamente si la excitaci\'on efectiva se encuentra
preferentemente en la rama polarit\'onica superior o en la inferior.

\subsection{Fidelidad del modelo efectivo}

La evoluci\'on obtenida en esta nota corresponde al hamiltoniano efectivo,
despu\'es del desplazamiento y de la aproximaci\'on de onda rotante
polarit\'on-fon\'on. Para cuantificar cu\'an bien reproduce la din\'amica del
hamiltoniano HC completo se comparan ambos estados mediante la fidelidad.
Escribiendo $\ket{\Psi_{\mathrm{HC}}(t)}$ para la evoluci\'on del hamiltoniano
completo y $\ket{\Psi_{\mathrm{ef}}(t)}$ para la evoluci\'on efectiva,
\begin{equation}
  \boxed{
  F(t)
  =
  \left|
  \braket{\Psi_{\mathrm{HC}}(t)|\Psi_{\mathrm{ef}}(t)}
  \right|^2.
  }
\end{equation}
Esta expresi\'on corresponde a la notaci\'on
$|\braket{\Psi(t)|\tilde\Psi(t)}|^2$ del manuscrito, pero las etiquetas
expl\'icitas evitan confundir el estado efectivo con el estado desplazado.
La fidelidad satisface
\begin{equation}
  0\leq F(t)\leq1.
\end{equation}
Si $F(t)\simeq1$, el hamiltoniano efectivo reproduce tanto las poblaciones como
las fases relativas relevantes del estado completo. Una fidelidad peque\~na
indica que los t\'erminos descartados han acumulado efectos importantes.

\section{Lectura f\'isica de la soluci\'on}

La soluci\'on exacta muestra que la din\'amica del modelo HC efectivo no es una
oscilaci\'on \'unica. Cada par $(n,m)$ tiene su propia frecuencia
\begin{equation}
  \Omega_{nm}^{\mathrm{ef}}
  =
  \sqrt{(2g\sqrt n-\nu)^2+\nu^2\lambda^2m}.
\end{equation}
La distribuci\'on fot\'onica selecciona varios valores de $n$ y la
distribuci\'on vibracional selecciona varios valores de $m$. El estado total
es la superposici\'on coherente de todas esas oscilaciones.

En resonancia polarit\'on-fon\'on, $\delta_n=0$, la frecuencia se reduce a
\begin{equation}
  \boxed{
  \Omega_{nm}^{\mathrm{ef}}
  =
  |k_m|
  =
  \nu|\lambda|\sqrt m,
  \qquad
  \delta_n=0.
  }
\end{equation}
En este caso la mezcla es m\'axima,
$C_{nm}=S_{nm}=1/\sqrt2$, y la transferencia entre
$\ket{+,n,m-1}$ y $\ket{-,n,m}$ puede ser completa.

Lejos de resonancia, $|\delta_n|\gg|k_m|$,
\begin{equation}
  \Omega_{nm}^{\mathrm{ef}}
  \simeq|\delta_n|,
  \qquad
  \frac{|k_m|}{\Omega_{nm}^{\mathrm{ef}}}\ll1.
\end{equation}
La amplitud de transferencia queda suprimida y los estados del doblete
principalmente acumulan fases distintas.

La aproximaci\'on coherente a\'nade una imagen geom\'etrica. Las dos ramas de
energ\'ia efectiva correlacionan los polaritones con dos paquetes
vibracionales que rotan en direcciones opuestas. Esta separaci\'on convierte a
la vibraci\'on en un registro de la historia polarit\'onica. Precisamente por
eso, al extender el modelo a dos mol\'eculas, la din\'amica vibracional puede
participar en la generaci\'on, almacenamiento o degradaci\'on de correlaciones
entre excitaciones moleculares mediadas por la cavidad.

\bigskip

La organizaci\'on por dobletes permite conservar una lectura simple incluso
cuando el estado inicial contiene muchas excitaciones fot\'onicas y
vibracionales: primero se determina qu\'e amplitud inicial entra a cada
doblete mediante $a_{nm}$ y $b_{nm}$; despu\'es cada doblete evoluciona con
$A_{nm}$, $\delta_n$, $k_m$ y $\Omega_{nm}^{\mathrm{ef}}$; finalmente, todas
las contribuciones se suman para reconstruir el estado completo y sus
observables.

\end{document}
